\documentclass[UTF8,a4paper,12pt]{ctexart}

\usepackage[margin=1in]{geometry}
\usepackage{array}
\usepackage{booktabs}
\usepackage{float}
\usepackage{xcolor}
\usepackage{hyperref}
\setlength{\emergencystretch}{3em}

\definecolor{reportlink}{RGB}{65, 145, 210}

\hypersetup{
  colorlinks=true,
  linkcolor=reportlink,
  urlcolor=reportlink,
  citecolor=reportlink
}

\newcommand{\project}{DiPECS}

\title{\project{} 中期报告}
\author{\project{} Team}
\date{2026 年 4 月 27 日}

\begin{document}

\maketitle

\begin{abstract}
中期阶段，\project{} 完成了从项目设想到最小可执行原型的转变。项目目标是构建一个面向 Android 的
用户意图预测与本地资源优化系统：在用户动作发生之前，通过本地采集和隐私脱敏形成结构化上下文，再由
本地规则或云端模型输出候选意图，最后由本地策略引擎决定是否执行低风险优化动作。到中期节点，项目已经
完成 Rust 核心类型体系、隐私空气隙、窗口聚合、规则式意图生成、策略校验、动作执行骨架、动作总线和
离线测试框架；Android 侧完成了 UsageStats、通知和基础系统状态等采集方向的原型验证。下一阶段重点是
打通 Android 到 Rust 的生产入口、完善 golden trace/replay、接入真实云端后端，并验证低风险动作是否带来
可测的系统收益。
\end{abstract}

\newpage
{\hypersetup{linkcolor=black}\tableofcontents}
\newpage

\section{项目定位}

\project{} 是一个面向 Android 端侧环境的意图预测系统原型。它的核心命题是：在用户动作实际发生之前，
系统能否根据近期上下文预测用户可能需要的资源，并提前进行低风险准备，从而降低冷启动或跨应用切换时的
可感知延迟。

中期汇报时，项目目标平台设定为 Android 13（API 33）和可定制 AOSP 环境；实现语言采用 Rust + Kotlin。
Rust 负责核心逻辑，包括事件 schema、隐私脱敏、上下文聚合、策略校验、动作状态机和离线回放；Kotlin 负责
Android 应用层，包括权限、服务、系统 API 调用和用户界面。

项目不把云端 LLM 视为直接控制设备的执行者。LLM 或本地规则只能输出结构化建议，例如 intent、confidence、
risk 和 suggested actions；本地 Rust 策略模块保留最终执行权。这是中期阶段最重要的系统边界。

\section{系统闭环}

中期阶段定义的闭环是：

\begin{center}
\small
本地采集 $\rightarrow$ 隐私脱敏 $\rightarrow$ 上下文聚合 $\rightarrow$ 意图预测
$\rightarrow$ 策略校验 $\rightarrow$ 本地优化执行 $\rightarrow$ 结果记录
\end{center}

各阶段含义如下：

\begin{table}[H]
\centering
\small
\begin{tabular}{p{0.20\linewidth}p{0.68\linewidth}}
\toprule
阶段 & 内容 \\
\midrule
本地采集 & 通过 Android 公开 API 和系统状态接口采集应用切换、通知、时间、网络、电量等信号。 \\
隐私脱敏 & 将通知文本、位置、文件路径等敏感信息转成粗粒度标签或摘要，避免 raw data 进入决策层。 \\
上下文聚合 & 将连续事件按时间窗口聚合为 StructuredContext，形成可输入模型或规则的统一上下文。 \\
意图预测 & 由规则基线或云端模型输出候选 intent、置信度、风险等级和建议动作。 \\
策略校验 & PolicyEngine 根据风险、置信度、能力等级和动作类型判断是否允许执行。 \\
本地执行 & 只执行低风险优化，如 KeepAlive、Prewarm/Prefetch 的安全子集或记录 hint。 \\
结果记录 & 通过 trace、audit record 和 golden replay 验证行为是否可复现、可审计。 \\
\bottomrule
\end{tabular}
\caption{中期阶段定义的系统闭环}
\end{table}

\section{平台选择}

中期阶段选择 Android 13 作为目标平台，主要基于稳定性和可验证性。

\begin{itemize}
  \item Android 13 的公开 API 相对成熟，UsageStatsManager、NotificationListenerService 和基础设备状态接口
  足以验证第一阶段算法闭环 \cite{android-usagestats,android-notification}。
  \item AOSP 环境便于后续从普通 APK 原型逐步下沉到 root、daemon 或 framework 层，测试更强的资源调度能力。
  \item Rust 交叉编译和 Android NDK 工具链已经可以支持核心库复用，降低后续 Kotlin/Rust 混合开发风险。
  \item Android 12 之后的系统具备 eBPF 基础能力，为后续内核级事件观测保留扩展空间 \cite{aosp-ebpf}。
\end{itemize}

因此，中期设计采用渐进路线：先在普通应用权限内验证采集、脱敏和策略闭环；再在更高权限环境中评估真实资源收益。

\section{架构分层}

中期阶段完成了对 Rust workspace 的分层梳理，目标是避免上层业务逻辑污染底层 schema 和策略边界。
当前工程中的主要模块职责如下：

\begin{table}[H]
\centering
\small
\begin{tabular}{p{0.20\linewidth}p{0.25\linewidth}p{0.45\linewidth}}
\toprule
层级 & 模块 & 职责 \\
\midrule
协议层 & aios-spec & 定义 RawEvent、SanitizedEvent、StructuredContext、Intent、Action 和 trace schema。 \\
核心层 & aios-core & 实现隐私空气隙、窗口聚合、策略引擎、动作生命周期和 trace 验证逻辑。 \\
采集层 & aios-collector & 接入 Android JSONL、本机系统采集、进程/系统状态读取和采集统计。 \\
决策层 & aios-agent & 提供 RuleBased、Cloud/Mock 后端和 DecisionRouter。 \\
动作层 & aios-action & 接收 AuthorizedAction，提供离线 adapter 和默认执行器骨架。 \\
工具层 & aios-cli / aios-daemon & 提供 replay、pipeline、daemon runtime 和观测入口。 \\
\bottomrule
\end{tabular}
\caption{中期阶段的工程分层}
\end{table}

这套分层体现了机制与策略分离。云端或 agent 层可以提出建议，但 schema、policy 和 action lifecycle 位于
本地核心层，负责最终边界控制。

\section{机制与策略分离}

中期汇报中强调了“云端负责 policy-like 推理，本地保留 mechanism 执行权”的原则。更准确地说，模型层负责
意图理解和候选动作生成；本地核心负责确定性审查和执行。

\begin{table}[H]
\centering
\small
\begin{tabular}{p{0.20\linewidth}p{0.34\linewidth}p{0.34\linewidth}}
\toprule
维度 & 云端/决策后端 & 本地 Rust core \\
\midrule
性质 & 非确定性推理或规则判断 & 确定性状态机和策略校验 \\
职责 & 意图理解、模式推断、候选动作生成 & 脱敏、上下文聚合、权限与风险校验、审计 \\
状态 & 可失败、可降级、可替换 & 可回放、可测试、可追责 \\
安全边界 & 只输出结构化建议 & 保留最终授权和执行权 \\
\bottomrule
\end{tabular}
\caption{机制与策略分离}
\end{table}

该原则避免了让 LLM 直接调用 Android API 或系统命令。即使模型输出了动作，本地也会根据风险等级、置信度、
动作白名单和上下文 target 检查决定是否授权。

\section{已完成工作}

\subsection{基础采集原型}

Android 侧已经围绕公开 API 验证基础采集方向：

\begin{itemize}
  \item 使用 UsageStatsManager 采集应用使用和前后台切换信息。
  \item 使用 NotificationListenerService 采集通知来源、类别、channel 和出现/撤销事件。
  \item 采集基础设备上下文，如时间、网络、电量和屏幕状态 \cite{android-battery,android-power}。
  \item 通过 append-only JSONL 或等价 trace 形式，将 Android 事件转入 Rust 侧处理链路。
\end{itemize}

这些采集源可以支撑第一阶段用户意图判断。中期阶段尚未把 Accessibility、Binder trace 或 eBPF 作为主路径，
而是保留为后续增强方向 \cite{android-accessibility}。

\subsection{Rust 核心链路}

中期阶段 Rust 侧已经完成最小闭环所需的主要模块：

\begin{itemize}
  \item \textbf{事件模型}：在 aios-spec 中定义 RawEvent、SanitizedEvent、StructuredContext、IntentBatch、
  SuggestedAction 和 trace 相关类型。
  \item \textbf{隐私脱敏}：PrivacyAirGap 将原始事件转换为脱敏事件，确保通知正文、文件路径等敏感字段不会进入后续后端。
  \item \textbf{窗口聚合}：WindowAggregator 将事件按时间窗口组织为 StructuredContext，并生成 ContextSummary。
  \item \textbf{意图生成}：规则后端/模拟后端可以根据通知、应用切换、屏幕状态、系统状态等信号生成候选 intent。
  \item \textbf{策略校验}：PolicyEngine 根据风险、置信度和动作类型过滤不安全建议。
  \item \textbf{动作执行骨架}：DefaultActionExecutor 和 OfflineAdapter 提供动作执行与模拟接口。
  \item \textbf{动作总线}：ActionBus 使用 channel 解耦采集、处理和执行链路。
\end{itemize}

\subsection{测试与验证}

中期阶段已经建立了覆盖核心模块的测试体系。测试重点包括：

\begin{itemize}
  \item 隐私脱敏：验证 raw PII 不进入 SanitizedEvent 或 StructuredContext。
  \item 窗口聚合：验证 foreground apps、notified apps、semantic hints、system status summary 的构建。
  \item 策略引擎：验证风险等级、置信度、动作白名单和 target 检查。
  \item 动作执行：验证不同 action type 的执行骨架和错误处理。
  \item 模拟决策：验证多类信号能被转为 intent/action。
  \item 动作总线：验证采集与处理之间的异步边界。
\end{itemize}

GoldenTrace 数据结构和离线 replay 方向已经确定。中期之后，项目需要把它从测试骨架推进为稳定的回归基线。

\section{为什么需要这个系统}

\subsection{移动端体验痛点}

移动端用户经常在聊天、浏览器、文档、开发工具或办公应用之间切换。传统系统通常在资源压力出现后才被动回收或
重建进程；当用户再次打开应用时，Dex 加载、资源解析、数据库初始化和后台状态恢复都可能造成可感知延迟。

\project{} 的思路是把部分冷启动成本提前：在高置信度场景中，系统可以保持短窗口内的关键进程或缓存资源，
将部分冷启动转为热启动。中期阶段的目标不是立刻证明所有场景都有收益，而是先建立能够安全提出和审查这类动作的
系统框架。FALCON 等移动端预启动工作说明，使用位置、时间等上下文预测 app 启动并做 prelaunch 是可研究的
系统优化方向 \cite{falcon}。

\subsection{语义鸿沟}

Android 能提供大量低层事件，例如包名、时间戳、UID、PID、通知 channel、网络和内存状态，但系统并不知道用户
为什么切换应用、某条通知是否重要、某个进程是否值得保留。换言之，系统知道 What/When，却缺少 Why 和
What-Inside。

\project{} 尝试在 PID 级事件和 intent 级语义之间建立中间层：用脱敏上下文和规则/模型生成候选意图，再用本地
策略决定动作。这是项目相对普通日志采集器的主要区别。

\section{技术路线}

\subsection{Android 信息来源}

中期设计把信息来源分为三层：

\begin{table}[H]
\centering
\small
\begin{tabular}{p{0.24\linewidth}p{0.62\linewidth}}
\toprule
层级 & 代表来源 \\
\midrule
公开 API & UsageStats、NotificationListener、Battery、Network、Screen state。 \\
系统特权/开发者接口 & Accessibility、ADB、dumpsys、Shizuku 或 root 环境。 \\
内核事件 & eBPF、Binder trace、文件 I/O、进程生命周期和内存压力事件。 \\
\bottomrule
\end{tabular}
\caption{中期规划中的 Android 信息来源}
\end{table}

中期原型主要落在第一层，确保普通授权环境下可运行；第二、三层作为后续 AOSP/root 环境验证方向。

\subsection{多源信息融合}

多源信息融合的核心是统一 schema。不同采集源输出格式不同、粒度不同、隐私风险不同，必须先归一化为 RawEvent，
再通过 PrivacyAirGap 进入 SanitizedEvent。窗口聚合之后，DecisionRouter 和 PolicyEngine 只面对统一的
StructuredContext，不需要理解每个采集源的原始细节。

这种设计使系统可以逐步扩展数据源。新增采集源只要能映射到 RawEvent 或 CollectorEnvelope，就能复用脱敏、
聚合、决策和审计链路。

\subsection{资源调度策略}

中期阶段定义了三类低风险动作方向：

\begin{table}[H]
\centering
\small
\begin{tabular}{p{0.20\linewidth}p{0.22\linewidth}p{0.46\linewidth}}
\toprule
动作方向 & 触发条件 & 内容 \\
\midrule
KeepAlive & 高置信度、低风险 & 短窗口保护相关进程或本系统任务，减少重建成本。 \\
Prewarm/Prefetch & 中高置信度、target 明确 & 预读可访问资源或预热本应用/安全资源。 \\
Hint/NoOp & 低置信度或风险不明 & 只记录候选，不做实际资源占用。 \\
\bottomrule
\end{tabular}
\caption{中期阶段定义的动作方向}
\end{table}

实际执行时，动作必须经过 policy 检查。中期阶段更关注动作生命周期和策略边界，而不是直接追求强系统控制能力。

\section{阶段性成果}

到中期节点，项目已经完成以下成果：

\begin{itemize}
  \item 明确了 Android 端侧 AIOS 原型的研究定位和系统边界。
  \item 建立了 Rust workspace 的协议层、核心层、采集层、决策层、动作层和工具层划分。
  \item 实现了 RawEvent 到 SanitizedEvent，再到 StructuredContext 的核心数据流。
  \item 建立了 PolicyEngine 和 ActionLifecycle 的早期形态，保证模型建议不会直接执行。
  \item 完成了动作执行骨架和离线 adapter，为 replay 和测试打下基础。
  \item 形成了基于测试和 golden trace 的可复现验证方向。
\end{itemize}

这些成果说明项目已经越过单纯方案设计阶段，进入可运行原型阶段。

\section{存在问题}

中期阶段仍存在以下问题：

\begin{itemize}
  \item Android 到 Rust 的生产入口仍需进一步稳定，JNI、JSONL 或 daemon 入口需要明确主路径。
  \item 云端 LLM 当时仍处于 mock/proxy 或接口设计阶段，真实云端后端需要后续接入。
  \item GoldenTrace 和 replay 还需要扩充数据集，覆盖通勤、学习、社交、办公等典型场景。
  \item 真实 Android action bridge 尚不完整，许多动作只能离线模拟或本地 stub。
  \item 内核级信息抽取、eBPF、Binder trace 和文件热点画像还停留在后续 root/AOSP 阶段。
\end{itemize}

\section{下一阶段计划}

中期之后的计划包括：

\begin{enumerate}
  \item 完成 Android 原型端到端演示：采集、脱敏、预测、策略校验和本地动作记录。
  \item 建立稳定 Golden Trace 数据集，并将 replay/audit hash 纳入回归测试。
  \item 对比无预热、启发式规则和模型预测三种策略的动作数量、拒绝率、延迟收益和资源开销。
  \item 接入真实云端 LLM 后端，同时保证网络失败时可以回退到 RuleBased 或 NoOp。
  \item 在权限更高的环境中验证 eBPF、Binder 和文件预读等增强能力，但不把它们作为普通 APK 阶段的前提。
\end{enumerate}

\section{核心贡献}

中期阶段形成的核心贡献可以概括为四点：

\begin{itemize}
  \item \textbf{隐私 Air-Gap 架构}：敏感信息在本地确定性脱敏，后续决策后端只看到结构化上下文。
  \item \textbf{确定性状态机和回放方向}：即使引入 LLM，系统仍能通过 trace 和 audit 记录复现关键行为。
  \item \textbf{多源信息融合框架}：将 API 层、系统状态和未来内核事件统一到同一事件 schema。
  \item \textbf{机制-策略分离}：模型提出建议，本地保留执行权，兼顾智能性和系统安全边界。
\end{itemize}

\section{结论}

中期阶段，\project{} 已经完成从概念到最小可执行原型的主要跨越。项目证明了在 Android 公开 API 和 Rust
核心链路基础上，可以建立一条可脱敏、可聚合、可决策、可审查、可回放的数据流。接下来的关键不再是单纯增加
更多规则，而是把 Android 真实输入、云端后端、golden trace 和 Android-safe 动作桥接接到同一条链路上，
并用实验指标验证预测与优化是否真正带来收益。

\begin{thebibliography}{99}

\bibitem{android-usagestats}
Android Developers.
\newblock UsageStatsManager.
\newblock \href{https://developer.android.com/reference/android/app/usage/UsageStatsManager}{Android API reference}.

\bibitem{android-notification}
Android Developers.
\newblock NotificationListenerService.
\newblock \href{https://developer.android.com/reference/android/service/notification/NotificationListenerService}{Android API reference}.

\bibitem{android-battery}
Android Developers.
\newblock BatteryManager.
\newblock \href{https://developer.android.com/reference/android/os/BatteryManager}{Android API reference}.

\bibitem{android-power}
Android Developers.
\newblock PowerManager.
\newblock \href{https://developer.android.com/reference/android/os/PowerManager}{Android API reference}.

\bibitem{android-accessibility}
Android Developers.
\newblock AccessibilityService.
\newblock \href{https://developer.android.com/reference/android/accessibilityservice/AccessibilityService}{Android API reference}.

\bibitem{aosp-ebpf}
Android Open Source Project.
\newblock Extend the kernel with eBPF.
\newblock \href{https://source.android.com/docs/core/architecture/kernel/bpf}{AOSP documentation}.

\bibitem{falcon}
Tingxin Yan et al.
\newblock FALCON: Fast App Launching via Predictive User Context.
\newblock \href{https://www.microsoft.com/en-us/research/publication/fast-app-launching-for-mobile-devices-using-predictive-user-context/}{Microsoft Research page}.

\end{thebibliography}

\end{document}
